\printedpage{4}
% head: A. Grothendieck : Dérivateurs, Ch. I, Édition des Universités de Montpellier II et Paris VII
% foot: Version 8 juillet 2001

\msp{page 4}

\section*{2. Heuristic preliminaries on derivators.}

We take a full subcategory $\Dia$ (as in ``diagrams'') of $\Cat$. It can be all of $\Cat$, or smaller. But we will most often assume that $\Dia$ contains at least the ``ordered categories'' (corresponding to ordered sets, \textit{i.e.} such that $a$) every morphism is a mono and $b$) every iso is an identity \texttt{[this is not true: the category reduced to a double arrow satisfies these conditions, but is not an ordered set]}). By abuse of language, we identify an ordered set with the corresponding category (and a category with the topos it defines?). We regard $\Dia$ as a 2-\textit{category}.\footnote{N.B. We need stability conditions for $\Dia$, at least stability under $I \times J$ and $I \amalg J$, in fact, considerably more \dots}

A \textit{derivator} $\DD$ (on $\Dia$, as the ``\textit{domain}'' of the derivator) is a contravariant 2-functor from $\Dia$ to $\Cat$, satisfying certain conditions. Thus
\[ \DD : \Dia^{\circ} \longrightarrow \Cat \quad . \]
If $I \in \Ob_0 \Dia$ is a category (or a ``type of diagram'') in the domain of $\DD$, we must therefore have a category $\DD(I)$. The objects of $\DD(I)$ will be called the \textit{coefficient-objects} (or \textit{coefficients}) of type $\DD$, over $I$. For
\[ f : I \longrightarrow J \]
a 1-arrow of $\Dia$, we must have a functor
\[ f^*_{\DD} \text{ or } f^* : \DD(J) \longrightarrow \DD(I) \quad , \]
called the \textit{``inverse image functor''}. If $F$ is a coefficient of type $\DD$ on $J$ its ``\textit{inverse image}'' is therefore a coefficient of type $\DD$ on $I$. We assume (to simplify!) that the 2-functor $\DD$ is \textit{strict}, hence that we have strict transitivity
\[ (gf)^* = f^*g^* \ , \qquad \text{if} \qquad I \xrightarrow{f} J \xrightarrow{g} K \qquad \text{in} \ \ \Dia \ . \]

\msp{page 5}

Moreover, if we have two functors $f, g : I \rightrightarrows J$, and a morphism between them
\[
\begin{tikzcd}[column sep=large]
I \arrow[r, shift left=2, "f"] \arrow[r, shift right=2, "g"'] \arrow[r, phantom, "{\downarrow}u"] & J
\end{tikzcd}
\quad ,
\]
we take a morphism $u^*_{\DD}$ or $u^*$ between the functors $g^*$ and $f^*$, whence
\[
\begin{tikzcd}[column sep=large]
\DD(I) & \DD(J) \arrow[l, shift right=2, "f^*"'] \arrow[l, shift left=2, "g^*"] \arrow[l, phantom, "{\uparrow}u^*"]
\end{tikzcd}
\quad .
\]

\printedpage{5}
% head: A. Grothendieck : Dérivateurs, Ch. I, Édition des Universités de Montpellier II et Paris VII
% foot: Version 8 juillet 2001

We must also have transitivity for a composition of morphisms of functors:
\[
\begin{tikzcd}[column sep=large]
I \arrow[r, shift left=4, "f"] \arrow[r, "g"' near start] \arrow[r, shift right=4, "h"'] \arrow[r, phantom, shift left=2, "{\downarrow}u"] \arrow[r, phantom, shift right=2, "{\downarrow}v"] & J
\end{tikzcd}
\]
must give
\[ (vu)^* = u^*v^* \]
thus
\[
\begin{tikzcd}[column sep=large]
\DD(I) & \DD(J) \arrow[l, shift right=4, "f^*"'] \arrow[l, "g^*" near end] \arrow[l, shift left=4, "h^*"] \arrow[l, phantom, shift right=2, "{\uparrow}u^*"] \arrow[l, phantom, shift left=2, "{\uparrow}v^*"]
\end{tikzcd}
\quad .
\]
Moreover, if $f$ is an identity, $f^*$ too
\[ (\id_I)^* = \id_{\DD(I)} \quad , \]
and likewise if $u$ is an identity
\[ (\id_f)^* = \id_{f^*} \quad . \]

\msp{page 6}

I indicate right away the most important axiom on a derivator (beyond the datum of a strict 2-contra functor):

\textbf{Der}$_1$: If $f : I \to J$ in $\Dia$, then $f^*_{\DD} = f^* : \DD(J) \to \DD(I)$ has a right adjoint and a left adjoint.

These functors will be denoted (at least provisionally) $f^{\DD}_*$ or $f_*$ , and $f^{\DD}_!$ or $f_!$ , that is
\[ f_! \, , \ f_* : \DD(I) \rightrightarrows \DD(J) \quad , \]
so that we find a sequence of \textit{three} adjoint functors between $\DD(I)$ and $\DD(J)$
\[ f_! \, , \ f^*, \ f_* \quad . \]

N.B. For most of the developments we will not need the existence of both the $f_!$ and the $f_*$ . A \textit{left derivator} requires only the $f_*$ (right adjoints of the $f^*$), and the existence of $f_!$ only when $f$ is an \textit{open immersion}. Dually, a \textit{right derivator} requires only the $f_!$ (left adjoints of the $f^*$), and $f_*$ only if $f$ is a \textit{closed immersion}. It is probably the left derivators that are the most important.

\msp{page 7}

\textbf{Der}$_2$: If $f : I \to J$ is in $\Dia$, if it is a \textit{closed} immersion, the right biadjoint $f^!_{\DD}$ or $f^!$ of $f^*$ exists. If it is an open immersion, the left biadjoint $f^?$ exists.\footnote{The biadjoints should be required only for a pointed derivator. \textsf{HOT} does not satisfy this.} In
\printedpage{6}
% head: A. Grothendieck : Dérivateurs, Ch. I, Édition des Universités de Montpellier II et Paris VII
% foot: Version 8 juillet 2001
both cases, we want $f^*$ to be a localization functor, which is equivalent to saying that $f_*$ is fully faithful, or that $f_!$ is fully faithful.\footnote{We must in addition characterize the image of $f_*$ if $f$ is a closed immersion (the $F$ such that $g^*F \overset{\sim}{\to} e_{\DD(U)}$) and that of $f_!$ if it is an open immersion (the $F$ such that $g^*F \simeq \emptyset_{\DD(Y)}$), where $g$ is the complementary immersion.}

N.B. In the case of a left derivator, we assume only the existence of $f^!$ if $f$ is a closed immersion. Thus to summarize, for such a derivator, we assume the existence of
\[
\left\{
\begin{aligned}
& f_* \text{ for any } f : I \to J \text{ in } \Dia \\
& f_! \text{ if } f \text{ is an open immersion} \\
& f^! \text{ if } f \text{ is a closed immersion (pointed case only).}
\end{aligned}
\right.
\]
Dually, for a right derivator, we assume the existence of
\[
\left\{
\begin{aligned}
& f_! \text{ for any } f : I \to J \text{ in } \Dia \\
& f_* \text{ if } f \text{ is a closed immersion} \\
& f^? \text{ if } f \text{ is an open immersion (pointed case only).}
\end{aligned}
\right.
\]
I would like to formulate clever axioms, for the case of a pair
\[ Y \overset{i}{\hookrightarrow} I \overset{j}{\hookleftarrow} U \]
of an open immersion $j$ and a closed immersion $i$ in $\Dia$, complementary to one another. First a formulation if we assume $\DD$ ``\textit{pointed}'', \textit{i.e.} the categories $\DD(I)$ pointed, and the functors $f^*$ compatible with the pointings. We want to ``interpret'' the square (1) on page 2:

\msp{page 8}

\[
\begin{tikzcd}
& F \arrow[dr, "\mathrm{adj}"] & \\
j_!j^*F \arrow[ur, "\mathrm{adj}"] \arrow[dr] & & i_*i^*F \\
& e \arrow[ur] &
\end{tikzcd}
\quad ,
\tag*{(1)}
\]

\textbf{Der}$'_2$: For all $I, J \in \Dia$, we have
\[ \DD(I \amalg J) \xrightarrow[(\alpha^*, \beta^*)]{\sim} \DD(I) \times \DD(J) \quad , \]
where $\alpha, \beta$ are the canonical injections of $I$ and $J$ into $I \amalg J$.

\printedpage{7}
% head: A. Grothendieck : Dérivateurs, Ch. I, Édition des Universités de Montpellier II et Paris VII
% foot: Version 8 juillet 2001

square both cartesian and cocartesian. We are thus looking for a square in $\DD(I)$, but better than that. Let $J = \Del_1 \times \Del_1$ be the standard square
\[
\begin{tikzcd}
& (1,0) \arrow[dl] & \\
(1,1) & & (0,0) \arrow[ul] \arrow[dl] \\
& (0,1) \arrow[ul] &
\end{tikzcd}
\]
and let us consider the category $I \times J$. We have
\[
\begin{aligned}
& J \longrightarrow \Hom(I, I \times J) \\
& j \longmapsto \bigl(i \longmapsto (i,j)\bigr)
\end{aligned}
\]
On the other hand, by $\DD$
\[ \Hom(I, I \times J)^{\circ} \longrightarrow \Hom(\DD(I \times J), \DD(I)) \]
whence a composite
\[ J^{\circ} \longrightarrow \Hom(\DD(I \times J), \DD(I)) \]
or again
\[ \DD(I \times J) \longrightarrow \Hom(J^{\circ}, \DD(I)) \quad . \]
This holds for every $J \in \Ob \, \Dia$. We can therefore say that the datum of a ``coefficient'' $\Phi$ on a product category $I \times J$ gives rise to a diagram of type $J^{\circ}$ in the category $\DD(I)$ of coefficients on $I$. The philosophy
\msp{page 9}
of derivators consists in large part, conversely, whenever we have a diagram in a category of coefficients $\DD(I)$ (on a diagram $I$), say diagram of type $K$: $K \to \DD(I)$, in asking ourselves whether it does not come ``canonically'' from an object of $\DD(I \times K^{\circ})$ \textit{i.e.} from a coefficient on $I \times K^{\circ}$.

Thus for the square p. 8. We would like it to come from $\DD(I \times J)$, where $J = \Del_1 \times \Del_1$, and this functorially, so we want a ``canonical'' functor
\[ \DD(I) \longrightarrow \DD(I \times \Del_1 \times \Del_1) \]
which on composing with
\[ \DD(I \times \Del_1 \times \Del_1) \longrightarrow \Hom((\Del_1 \times \Del_1)^{\circ}, \DD(I)) \]
\printedpage{8}
% head: A. Grothendieck : Dérivateurs, Ch. I, Édition des Universités de Montpellier II et Paris VII
% foot: Version 8 juillet 2001
gives the functor described by square (1). Moreover, we must express in the present context the ``bicartesian'' condition on square (1), in the context of page 7. Analogous question for square (2) on page 2:
\[
\begin{tikzcd}
& F \arrow[dr] & \\
i_*i^!F \arrow[ur] \arrow[dr] & & j_*j^*F \\
& e \arrow[ur] &
\end{tikzcd}
\quad ,
\tag*{(2)}
\]
where we must in addition interpret the cartesian condition. We are in fact going to define a coefficient on the following open subcategory $W$ of $\Del_1 \times \Del_1$
\[
\begin{tikzcd}
(1,0) & \\
& (0,0) \arrow[ul] \arrow[dl] \\
(0,1) &
\end{tikzcd}
\quad ,
\]

\msp{page 10}

more precisely on $W \times I$, and if $\alpha$
\[ \alpha : W \times I \longrightarrow \Del_1 \times \Del_1 \times I \]
is the canonical inclusion, we will define the ``cartesian'' square as the image under $\alpha_*$ of this coefficient. This thus brings us back to the problem of first ``globalizing'' the partial diagram
\[
\begin{tikzcd}
F \arrow[dr] & \\
& j_*j^*F \\
e \arrow[ur] &
\end{tikzcd}
\quad .
\]
And likewise for square (1) except that we are at a loss if we choose to look at it mainly as cartesian (which suggests proceeding as just now) or as \underline{co}cartesian (which % verified transcription correction: the edition prints embarrassé, with double r